% =============================================================================
\section{Network Graph Representation}
\label{sec:network}
% =============================================================================

The multi-tank system is described as a directed graph
\begin{equation}
  G = (\Nodes,\, \Edges),
\end{equation}
where $\Nodes$ is the set of \emph{nodes} and $\Edges \subseteq \Nodes\times\Nodes$
is the set of directed \emph{edges}.  Every mass-transfer path in the system is
represented by an edge; the governing equations (Section~\ref{sec:dae}) are then
assembled purely from graph incidence.

% -----------------------------------------------------------------------------
\subsection{Nodes}
\label{sec:network:nodes}
% -----------------------------------------------------------------------------

Three node types are recognised:

\begin{center}
\begin{tabularx}{\linewidth}{@{} l l X @{}}
\toprule
\textbf{Type} & \textbf{YAML \texttt{type}} & \textbf{Description} \\
\midrule
\texttt{tank}       & \texttt{tank}       &
  Finite-volume hydrogen storage vessel with evolving thermodynamic state
  $(m_i, T_i, T_{s,i})$. \\[4pt]
\texttt{source}     & \texttt{source}     &
  Infinite reservoir that supplies mass at a prescribed enthalpy (e.g.\ a
  cryogenic pump feeding a refuel edge). \\[4pt]
\texttt{sink}       & \texttt{sink}       &
  Infinite reservoir that absorbs mass with no back-pressure (e.g.\
  \texttt{atmosphere} for venting, \texttt{fuel\_cell} for mission discharge). \\[4pt]
\texttt{junction}   & \texttt{junction}   &
  Stateless flow-splitting or merging point.  Conserves mass and enthalpy
  across its incident edges; split fractions are declared on the outflow
  edges.  Carries no dynamic state. \\
\bottomrule
\end{tabularx}
\end{center}

Only \texttt{tank} nodes carry dynamic states; \texttt{source}, \texttt{sink},
and \texttt{junction} nodes act as boundary conditions or purely algebraic
routing elements.

% -----------------------------------------------------------------------------
\subsection{Edges}
\label{sec:network:edges}
% -----------------------------------------------------------------------------

Four edge types are currently implemented, summarised in
\cref{tab:edge-types}.

\begin{table}[ht]
\centering
\caption{Edge types, their activation conditions, and the enthalpy assigned to
the transferred mass.}
\label{tab:edge-types}
\begin{tabularx}{\linewidth}{@{} l l X l @{}}
\toprule
\textbf{Type} & \textbf{Direction} & \textbf{Activation / flow rate} & \textbf{Enthalpy} \\
\midrule
\texttt{coupling}  & tank $\to$ tank &
  Activated by a valve model (e.g.\ pressure-triggered, PID, governor)
  when the pressure differential across the edge meets the opening criterion. &
  $h_\text{source}(T_\text{src},\rho_\text{src})$ \\[6pt]

\texttt{discharge} & tank $\to$ sink &
  Flow rate prescribed by a mission CSV file, $\mdot_\text{CSV}(t)$.
  In Config~B the rate continues at the mission value; a compensating
  heat flux is supplied to maintain the pressure floor
  (see \cref{sec:dae:configB}). &
  $h(T_i,\rho_i)$ \\[6pt]

\texttt{vent}      & tank $\to$ sink &
  Activated when $P_i \ge P_{\text{vent},i}$.
  Proportional rate model: $\mdot_\text{vent} = k_v\max(0, P_i - P_{\text{vent},i})$. &
  $h(T_i,\rho_i)$ \\[6pt]

\texttt{refuel}    & source $\to$ tank &
  Flow rate prescribed by a mission CSV or fixed schedule.
  Enthalpy computed from cryogenic pump isentropic model. &
  $h_\text{cryo}(P_i)$ \\
\bottomrule
\end{tabularx}
\end{table}

% -----------------------------------------------------------------------------
\subsection{Junction Nodes}
\label{sec:network:junctions}
% -----------------------------------------------------------------------------

A junction node is a stateless routing element that splits or merges mass
streams.  It carries no dynamic states and requires no thermodynamic model;
all computations are purely algebraic at each RHS evaluation.

\subsubsection{Splitter (one inflow, multiple outflows)}

Let $e_0$ be the single inflow edge and $e_1,\dots,e_K$ the $K$ outflow edges.
Each outflow edge declares a dimensionless \texttt{split\_fraction} $f_k \in
[0,1]$, with $\sum_{k=1}^{K} f_k = 1$.  The outflow rates are
\begin{equation}
  \mdot_{e_k} = f_k\,\mdot_{e_0}, \qquad k = 1,\dots,K.
  \label{eq:junction-split}
\end{equation}
Specific enthalpy is conserved at an ideal junction (no heat addition or shaft
work):
\begin{equation}
  h_{e_k} = h_{e_0}, \qquad k = 1,\dots,K.
  \label{eq:junction-split-h}
\end{equation}

\subsubsection{Merger (multiple inflows, one outflow)}

Let $e_1,\dots,e_J$ be inflow edges and $e_0$ the single outflow edge.
Mass conservation gives
\begin{equation}
  \mdot_{e_0} = \sum_{j=1}^{J} \mdot_{e_j}.
  \label{eq:junction-merge}
\end{equation}
The mixture enthalpy at the junction outlet is the mass-flow-weighted average:
\begin{equation}
  h_{e_0} = \frac{\displaystyle\sum_{j=1}^{J} \mdot_{e_j}\,h_{e_j}}{\mdot_{e_0}}.
  \label{eq:junction-merge-h}
\end{equation}

% -----------------------------------------------------------------------------
\subsection{Peripheral Components on Edges}
\label{sec:network:peripheral}
% -----------------------------------------------------------------------------

A \emph{peripheral component chain} can be attached to any edge to condition
the fluid stream between the source and target nodes.  The chain preserves the
mass flow rate (steady-state assumption) but transforms the thermodynamic state
$(P, T, h)$ of the stream.  The outlet enthalpy $h_\text{out}$ from the last
component in the chain is substituted for $h_\text{src}$ when assembling the
downstream energy balance \cref{eq:energy-balance}.

\subsubsection{Stream transformation model}

Let the raw stream at the edge source be $s_0 = (P_\text{src},\; T_\text{src},\;
\mdot_e,\; h_\text{src})$.  Each component $c_k$ maps
\begin{equation}
  s_k = c_k(s_{k-1}), \qquad k = 1,\dots,K,
\end{equation}
and the outlet state $s_K = (P_\text{out},\; T_\text{out},\; \mdot_e,\;
h_\text{out})$ is used in the energy balance at the downstream node.

\subsubsection{Implemented component models}

\paragraph{Ideal heat exchanger (\texttt{ideal\_heat\_exchanger}).}
Sets the stream to a prescribed target temperature $T_\text{tgt}$ with an
optional pressure drop $\Delta P$:
\begin{equation}
  P_\text{out} = P_\text{in} - \Delta P,\qquad
  T_\text{out} = T_\text{tgt},\qquad
  h_\text{out} = h(P_\text{out},\, T_\text{tgt}).
\end{equation}

\paragraph{Compressor (\texttt{compressor}).}
Raises stream pressure from $P_\text{in}$ to a specified $P_\text{out}$ with
isentropic efficiency $\eta$:
\begin{equation}
  h_\text{out} = h_\text{in} + \frac{h_{2s} - h_\text{in}}{\eta},\quad
  h_{2s} = h(P_\text{out},\, s_\text{in}).
  \label{eq:compressor}
\end{equation}
The outlet temperature and density are recovered from $(P_\text{out},
h_\text{out})$ via CoolProp.

\paragraph{Cryopump (\texttt{cryopump}).}
Pumps liquid hydrogen from a low-pressure reservoir at $P_0 = \SI{3}{\bar}$
to the downstream tank pressure $P_i$ with isentropic efficiency
$\eta_p = 0.78$.  Liquid saturation conditions ($x = 0$) are assumed at
suction:
\begin{equation}
  h_\text{out}(P_i)
  = h(P_0,\,x{=}0)
    + \frac{h_{2s}(P_i) - h(P_0,\,x{=}0)}{\eta_p},\quad
  h_{2s}(P_i) = h\!\left(P_i,\, s(P_0,x{=}0)\right).
  \label{eq:cryopump}
\end{equation}
Used on \texttt{refuel} edges.

% -----------------------------------------------------------------------------
\subsection{YAML Schema Mapping}
\label{sec:network:yaml}
% -----------------------------------------------------------------------------

The graph is declared in the \texttt{network} block of the scenario YAML file:

\begin{lstlisting}[language=Python, caption={Excerpt of YAML network declaration showing all node and edge types.}]
network:
  nodes:
    - node_id: lh2_tank
      type: tank          # tank | source | sink | junction
      fluid: LH2
      ...
    - node_id: gh2_buffer
      type: tank
      fluid: GH2
      ...
    - node_id: cryo_reservoir
      type: source        # infinite LH2 supply for refuelling
    - node_id: fuel_cell
      type: sink          # mission discharge endpoint
    - node_id: atmosphere
      type: sink          # venting endpoint
    - node_id: discharge_splitter
      type: junction      # splits flow between tank and fuel cell

  edges:
    # Coupling: LH2 tank feeds GH2 buffer via compressor + HEX chain
    - edge_id: lh2_to_buffer
      connection_type: coupling
      from_node: lh2_tank
      to_node: gh2_buffer
      components:
        - type: compressor
          parameters: {outlet_pressure: 35e5, efficiency: 0.82}
        - type: ideal_heat_exchanger
          parameters: {target_temperature: 293.0, pressure_drop: 0.5e5}

    # Discharge split: 95% to fuel cell, 5% to buffer top-up
    - edge_id: to_junction
      connection_type: discharge
      from_node: gh2_buffer
      to_node: discharge_splitter
    - edge_id: junction_to_fc
      connection_type: discharge
      from_node: discharge_splitter
      to_node: fuel_cell
      split_fraction: 0.95
    - edge_id: junction_to_buffer
      connection_type: coupling
      from_node: discharge_splitter
      to_node: lh2_tank
      split_fraction: 0.05

    # Vent edges
    - edge_id: vent_lh2
      connection_type: vent
      from_node: lh2_tank
      to_node: atmosphere
    - edge_id: vent_buffer
      connection_type: vent
      from_node: gh2_buffer
      to_node: atmosphere

    # Refuel edge with cryopump
    - edge_id: refuel_lh2
      connection_type: refuel
      from_node: cryo_reservoir
      to_node: lh2_tank
      components:
        - type: cryopump
          parameters: {reservoir_pressure: 3.0e5, efficiency: 0.78}
\end{lstlisting}

Mission discharge and venting edges that were previously implicit (hard-coded in
the dynamic model) are now explicit edges in the \texttt{network.edges} list.
Peripheral components are declared inline on the edge via the \texttt{components}
list; split fractions on junction outflow edges must sum to one.

% -----------------------------------------------------------------------------
\subsection{Example: CH2/CCH2 Two-Tank System}
\label{sec:network:example}
% -----------------------------------------------------------------------------

\Cref{fig:network-example} shows the directed graph for the reference scenario:
a 800-bar compressed hydrogen tank (CH2) coupled to a 400-bar cryo-compressed
hydrogen tank (CCH2) with mission discharge assigned to Tank~2.

\begin{figure}[ht]
\centering
\begin{tikzpicture}[
    node distance = 2.2cm and 3.0cm,
    tank/.style    = {draw, rounded corners=4pt, fill=blue!8, minimum width=2.8cm,
                      minimum height=1.2cm, align=center, font=\small\bfseries},
    boundary/.style= {draw, dashed, rounded corners=4pt, fill=gray!10,
                      minimum width=2.0cm, minimum height=0.9cm, align=center,
                      font=\small\itshape},
    edge/.style    = {-{Latex[length=6pt]}, thick},
    label/.style   = {font=\scriptsize, midway, align=center, text width=2.2cm},
  ]

  % Nodes
  \node[tank]     (t1)                      {Tank 1\\CH2\\800 bar};
  \node[tank]     (t2) [right=4.5cm of t1]  {Tank 2\\CCH2\\400 bar};
  \node[boundary] (fc) [right=3.5cm of t2]  {Fuel Cell\\(sink)};
  \node[boundary] (atm1) [above=1.8cm of t1]{Atmosphere\\(sink)};
  \node[boundary] (atm2) [above=1.8cm of t2]{Atmosphere\\(sink)};

  % Coupling edge
  \draw[edge, blue!70!black]
        (t1) -- node[label, above] {coupling\\$P_1 > P_2 + \Delta P$} (t2);

  % Vent edges
  \draw[edge, red!70!black]
        (t1) -- node[label, right] {vent\\$P_1 \ge P_{\text{vent},1}$} (atm1);
  \draw[edge, red!70!black]
        (t2) -- node[label, right] {vent\\$P_2 \ge P_{\text{vent},2}$} (atm2);

  % Discharge edge
  \draw[edge, green!50!black]
        (t2) -- node[label, above] {discharge\\$\mdot_\text{CSV}(t)$} (fc);

\end{tikzpicture}
\caption{Directed network graph for the CH2/CCH2 coupled system.  Tank~1 has no
discharge edge; Tank~2 carries the full mission discharge load.  Both tanks have
vent edges to the atmosphere sink.}
\label{fig:network-example}
\end{figure}
